% =====================================================================
%  CHAPITRE D'EXEMPLE -- démonstration des styles du polycopié
%  (à supprimer ou renommer une fois les styles adoptés)
% =====================================================================
\chapter{Guide des styles du polycopié}
\label{chap:styles}

Ce chapitre d'exemple illustre les conventions typographiques utilisées
dans l'ensemble du document : environnements \textit{exemple} (bleu) et
\textit{remarque importante} (rouge), tableaux, figures TikZ et figures
vectorielles exportées depuis Inkscape. Les grandeurs sont exprimées
avec \SI{}{} du paquet \texttt{siunitx} : \SI{48}{V}, \SI{10}{kW},
\SI{100}{kHz}.

% ----------------------------------------------------------------------
\section{Figures externes (SVG via Inkscape)}
% ----------------------------------------------------------------------

La figure~\ref{fig:convertisseur_buck} montre un hacheur série
abaisseur (\emph{buck}). Les figures sources sont des fichiers SVG
placés dans le dossier \texttt{Figures/} et compilés automatiquement
par le paquet \texttt{svg}.

\begin{figure}[htbp]
    \centering
    \includesvg[width=0.6\textwidth]{Industrial_drive_example.svg}
    \caption{Hacheur abaisseur de tension (\emph{buck}) : transistor,
      diode, inductance et condensateur de sortie.}
    \label{fig:convertisseur_buck}
\end{figure}

% ----------------------------------------------------------------------
\section{Figures TikZ et pgfplots}
% ----------------------------------------------------------------------

Les figures natives TikZ utilisent la palette institutionnelle.
La figure~\ref{fig:ondulation_courant} représente l'ondulation du
courant d'inductance en conduction continue.

\begin{figure}[htbp]
\centering
\begin{tikzpicture}
\begin{axis}[
  width=0.85\textwidth, height=5.5cm,
  xlabel={temps (\si{\milli\second})},
  ylabel={courant $i_L$ (\si{\ampere})},
  grid=major, grid style={black!12},
  axis line style={black!40},
  tick label style={font=\small},
  xmin=0, xmax=1, ymin=0, ymax=2.4,
]
  % courant triangulaire entre 1 et 2 A
  \addplot[accent, thick] coordinates
    {(0,1)(0.125,2)(0.375,1)(0.625,2)(0.875,1)(1,1.5)};
  % valeur moyenne
  \addplot[accenthesso, thick, dashed] coordinates {(0,1.5)(1,1.5)};
  % enveloppe de l'ondulation
  \addplot[textgrey, dotted, thick] coordinates {(0,2)(1,2)};
  \addplot[textgrey, dotted, thick] coordinates {(0,1)(1,1)};
  \legend{$i_L(t)$, $I_L = \SI{1.5}{A}$, $\Delta I_L$}
\end{axis}
\end{tikzpicture}
\caption{Ondulation du courant d'inductance $\Delta I_L = \SI{1}{A}$
  autour de sa valeur moyenne.}
\label{fig:ondulation_courant}
\end{figure}

% ----------------------------------------------------------------------
\section{Tableaux}
% ----------------------------------------------------------------------

Le tableau~\ref{tab:convertisseurs_cc} compare les principales
structures de convertisseurs DC/DC. Les tableaux utilisent
\texttt{booktabs} : pas de filets verticaux, uniquement
\texttt{toprule}, \texttt{midrule} et \texttt{bottomrule}.

\begin{table}[htbp]
\centering
\caption{Comparaison des convertisseurs DC/DC non isolés.}
\label{tab:convertisseurs_cc}
\small
\begin{tabular}{lccc}
\toprule
\textcolor{accenthesso}{\textbf{Structure}} &
\textcolor{accenthesso}{\textbf{Rapport}} &
\textcolor{accenthesso}{\textbf{Rendement}} &
\textcolor{accenthesso}{\textbf{Usage typique}} \\
\midrule
Buck (abaisseur)   & $V_o / V_e$           & élevé & alimentations, bus DC \\
Boost (élévateur)   & $V_e / (V_e - V_o)^{-1}$ & élevé & PFC, récupération \\
Buck-boost          & $-D/(1-D)$           & moyen & inversion, variation large \\
\'Cuk               & $-D/(1-D)$           & moyen & entrée/sortie en continu \\
\bottomrule
\end{tabular}
\end{table}

Note l'alignement des unités dans le texte courant :
$V_e = \SI{400}{V}$, $L = \SI{250}{\micro\henry}$.

% ----------------------------------------------------------------------
\section{Exemples (style bleu)}
% ----------------------------------------------------------------------

Les exemples apparaissent en bleu, avec filet supérieur et marges
réduites de part et d'autre.

\begin{example}[Calcul de l'ondulation de courant]
  Soit un hacheur buck alimentant une charge résistive sous
  $V_e = \SI{48}{V}$, $V_o = \SI{24}{V}$, avec une fréquence de
  découpage $f_s = \SI{100}{kHz}$ et une inductance
  $L = \SI{250}{\micro\henry}$.

  L'ondulation du courant d'inductance vaut
  \begin{equation}
    \Delta I_L = \frac{(V_e - V_o)\,D}{L\,f_s}
               = \frac{(48 - 24)\cdot 0{,}5}{250\mathrm{u}\cdot 10^5}
               \approx \SI{480}{mA}.
  \end{equation}
  Ce dimensionnement garantit la conduction continue dès que le
  courant moyen dépasse $\Delta I_L / 2 = \SI{240}{mA}$.
\end{example}

% ----------------------------------------------------------------------
\section{Remarques importantes (style rouge)}
% ----------------------------------------------------------------------

Les remarques critiques pour la sécurité ou le bon fonctionnement du
montage sont signalées en rouge.

\begin{important}[Ne jamais ouvrir le circuit d'inductance]
  L'interruption brutale du courant dans une inductance provoqu une
  surtension $v_L = -L\,\mathrm{d}i/\mathrm{d}t$ susceptible de
  détruire le transistor de découpage. Toujours vérifier la présence
  d'une diode de roue libre (ou d'un bras actif) \textbf{avant} de
  mettre sous tension le convertisseur.
\end{important}

\begin{important}[Choix du condensateur de sortie]
  Pour respecter une ondulation de tension de sortie
  $\Delta V_o < \SI{50}{mV}$, le condensateur doit présenter une
  résistance série (ESR) faible \textbf{et} une valeur suffisante :
  la composante ESR domine souvent à $f_s = \SI{100}{kHz}$.
\end{important}

% ----------------------------------------------------------------------
\section{Palette de couleurs du document}
% ----------------------------------------------------------------------

Le tableau~\ref{tab:palette} récapitule la palette utilisée dans tout
le polycopié : le rouge HEIG-VD est réservé aux remarques importantes
et aux alertes, le bleu HES-SO aux exemples et liens, les teintes
complémentaires aux graphiques.

\begin{table}[htbp]
\centering
\caption{Palette de couleurs du polycopié.}
\label{tab:palette}
\small
\begin{tabular}{llc}
\toprule
\textcolor{accenthesso}{\textbf{Couleur}} &
\textcolor{accenthesso}{\textbf{Usage}} &
\textcolor{accenthesso}{\textbf{RVB}} \\
\midrule
\textcolor{accent}{\rule{0.4cm}{0.25cm}} \texttt{accent} (HEIG-VD)
  & remarques importantes, alertes & 225, 37, 27 \\
\textcolor{accenthesso}{\rule{0.4cm}{0.25cm}} \texttt{accenthesso}
  & exemples, titres, liens & 0, 96, 156 \\
\textcolor{textgrey}{\rule{0.4cm}{0.25cm}} \texttt{textgrey}
  & texte secondaire, légendes & 150, 139, 131 \\
\midrule
\textcolor{accentgreen}{\rule{0.4cm}{0.25cm}} \texttt{accentgreen}
  & graphiques (série 1) & 0, 132, 86 \\
\textcolor{accentpurple}{\rule{0.4cm}{0.25cm}} \texttt{accentpurple}
  & graphiques (série 2) & 108, 60, 158 \\
\textcolor{accentorange}{\rule{0.4cm}{0.25cm}} \texttt{accentorange}
  & graphiques (série 3) & 232, 116, 17 \\
\textcolor{accentpink}{\rule{0.4cm}{0.25cm}} \texttt{accentpink}
  & graphiques (série 4) & 206, 54, 104 \\
\textcolor{accentteal}{\rule{0.4cm}{0.25cm}} \texttt{accentteal}
  & graphiques (série 5) & 0, 142, 148 \\
\textcolor{accentgold}{\rule{0.4cm}{0.25cm}} \texttt{accentgold}
  & graphiques (série 6) & 172, 134, 0 \\
\bottomrule
\end{tabular}
\end{table}

% ----------------------------------------------------------------------
\section{Listes et énumérations}
% ----------------------------------------------------------------------

Les listes à puces utilisent l'espacement compact défini globalement
(\texttt{itemsep=0pt}, marge droite réduite à \(2{,}5\)~em). Les
niveaux s'imbriquent avec des marqueurs distincts :

\begin{itemize}
  \item Premier niveau : puce ronde ;
  \item les items longs reviennent à la ligne avec la marge droite
    réduite, ce qui distingue visuellement le bloc liste du texte
    courant et évite que les lignes touchent la reliure ;
  \item imbrication sur trois niveaux :
  \begin{itemize}
    \item Deuxième niveau : tiret ;
    \item les sous-niveaux héritent de la même marge réduite ;
    \begin{itemize}
      \item Troisième niveau : point ;
      \item au-delà, préférer restructurer le contenu !
    \end{itemize}
  \end{itemize}
  \item retour au premier niveau.
\end{itemize}

Les énumérations numérotées ont leurs numéros en rouge HEIG-VD
(défini par \texttt{label=\textcolor{accent}{\textbackslash arabic*.}}
dans le préambule) :

\begin{enumerate}
  \item Dimensionner l'ondulation de courant $\Delta I_L$ ;
  \item vérifier la condition de conduction continue :
  $I_{load} > \Delta I_L / 2$ ;
  \item choisir la fréquence de découpage selon le compromis
    pertes--volume :
  \begin{enumerate}
    \item pertes de commutation croissantes avec $f_s$ ;
    \item volume des composants magnétiques décroissant avec $f_s$.
  \end{enumerate}
  \item valider thermiquement le choix des semi-conducteurs.
\end{enumerate}

% --- Variantes nommées : compact / aéré ---------------------------------
% (à définir une fois dans le préambule, cf. note en fin de section)

\begin{packedenum}
  \item liste compacte : \texttt{itemsep=0pt} ;
  \item utile pour les checklists, algorithmes, énumérations denses.
\end{packedenum}

\begin{airenum}
  \item liste aérée : \texttt{itemsep=1.5ex} ;
  \item utile pour les étapes d'un raisonnement long,
  \item ou les instructions de manipulation de laboratoire.
\end{airenum}

Les descriptions associatives restent utiles pour les glossaires de
symboles :

\begin{description}[font=\normalfont\bfseries\color{accenthesso}]
  \item[$D$] rapport cyclique du hacheur, $0 \le D \le 1$ ;
  \item[$f_s$] fréquence de découpage (\si{\kilo\hertz}) ;
  \item[$\eta$] rendement du convertisseur (\si{\percent}).
\end{description}

Enfin, les énumérations \emph{en ligne} (dans une phrase) s'obtiennent
avec la version étoilée du paquet \texttt{enumitem} : les étapes du
calcul sont a) identifier les intervalles de conduction, b) appliquer
la loi des mailles sur chaque intervalle, c) moyenner les équations
pour obtenir la relation entrée-sortie. (Voir note ci-dessous pour la
syntaxe exacte.)

% ----------------------------------------------------------------------
\section{Références croisées et citations}
% ----------------------------------------------------------------------

Le lecteur est renvoyé au chapitre~\ref{chap:styles}, à la
figure~\ref{fig:ondulation_courant} et au
tableau~\ref{tab:convertisseurs_cc}. Les références
bibliographiques suivent le style \texttt{plain}~\cite{mohan}.

% ============================================================
% RÉFÉRENCES
% ============================================================
\begin{thebibliography}{9}

\bibitem{teodorescu2011}
R.~Teodorescu, M.~Liserre, P.~Rodriguez,
\emph{Grid Converters for Photovoltaic and Wind Power Systems},
John Wiley \& Sons, 2011.
[Illustrations de référence : phaseurs spatiaux, réseaux forts/faibles,
défauts asymétriques.]

\bibitem{akagi}
H.~Akagi, E.~H.~Watanabe, M.~Aredes,
\emph{Instantaneous Power Theory and Applications to Power
Conditioning}, IEEE Press / Wiley, 2007.
[Théorie $p$-$q$ et filtres actifs.]

\bibitem{bossoney}
L.~Bossoney, \emph{Électronique de Puissance I, Partie 4 : Power
Factor Correction (PFC)}, HEIG-VD, 2008.
[Base des sections PFC et des exercices du chapitre 2.]

\bibitem{ieciec61000}
IEC 61000-2-4, \emph{Compatibility levels in industrial plants for
low-frequency conducted disturbances} ; IEC 61000-3-2, \emph{Limits
for harmonic current emissions}.
[Normes de compatibilité et d'émissions harmoniques.]

\bibitem{erickson}
R.~W.~Erickson, D.~Maksimović,
\emph{Fundamentals of Power Electronics}, 3\ieme{} éd., Springer, 2020.
[Convertisseurs DC/DC : boost, SEPIC, modes de conduction.]

\bibitem{mohan}
N.~Mohan, T.~M.~Undeland, W.~P.~Robbins,
\emph{Power Electronics: Converters, Applications, and Design},
3\ieme{} éd., Wiley, 2003.
[Applications industrielles, redresseurs, entraînements.]

\end{thebibliography}